\centerline{\bf \huge 7 класс}
\centerline{\normalsize{\it Работа рассчитана на \bf{180} минут}}

\problem{7.1}
В шестизначном числе пронумеровали слева направо все цифры числами от 1 до 6 , после чего все цифры с чётными номерами в том же порядке переставили на первые три места, а все цифры с нечётными номерами в том же порядке переставили на последние три места. Получившееся число совпало с исходным. Какое наибольшее количество различных цифр может быть в таком числе?

\problem{7.2}
Найдите наименьшее число, начинающееся с цифр 1234567 и делящееся на 225.

\problem{7.3}
В банке можно положить вклад на месяц под фиксированный процент. В конце месяца банк начислит проценты, а потом округлит (по обычным правилам) начисленную сумму до целого числа рублей. У Васи и Пети одинаковые суммы денег. Вася разделил свои деньги на две равные части и открыл два вклада, а Петя положил все свои деньги на один вклад. Может ли оказаться, что после выплаты процентов сумма у Васи будет меньше суммы у Пети?

\problem{7.4}
За круглый стол сели семь математиков. Каждому из них дали карточку, на которой написано число «1» или « $-1 »$. Затем на доску записали семь чисел - произведения чисел на карточках каждых двух математиков, сидящих рядом. После этого на доску записали восьмое число - произведение всех семи чисел на карточках математиков. Какое наименьшее значение может принимать сумма 8 выписанных на доске чисел?

\problem{7.5}
На острове живут рыцари и лжецы, рыцари всегда говорят правду, а лжецы всегда лгут. На чаепитие в зале собралось 30 человек. В процессе беседы десять из них сказали, что в зале ровно 20 рыцарей. Пятеро сказали, что в зале не больше 7 рыцарей. Двенадцать человек сказали, что в зале ровно 10 рыцарей. А остальные трое сказали, что в зале ровно 9 рыцарей. Каждый произнес ровно одну фразу. Сколько рыцарей присутствует на чаепитии?